\subsection{Step 1: Data Audit and Preprocessing}
\label{sec:step1_data_audit}

Step~1 converts the raw extracted JSON into a validated occurrence-level dataset that can support the later harmonization and graph-building stages. The step is intentionally non-destructive: it standardizes formatting, validates schema and vectors, expands list-like curricular provenance into occurrence rows, and records every repair or ambiguity in explicit audit files.

The expected raw fields are \texttt{skill\_id}, \texttt{name}, \texttt{description}, \texttt{sections}, \texttt{programmes}, \texttt{year}, \texttt{edu\_type}, \texttt{chunk\_id}, \texttt{name\_vector}, and \texttt{description\_vector}. The implemented preprocessing verifies that these fields are present, parses the two embedding columns as numeric arrays, checks dimensional consistency, normalizes text and metadata conservatively, and expands the raw records to one row per curricular occurrence.

The key correction introduced in this pass concerns curricular provenance. Direct broadcasting of \texttt{sections}, \texttt{programmes}, and \texttt{year} is attempted first. When that broadcasting is not structurally sufficient, the script falls back to the \texttt{chunk\_id}-encoded section-programme pairs together with a stable year lookup. These fallback repairs are not silently accepted: they are flagged with \texttt{provenance\_resolution\_confidence} and \texttt{requires\_manual\_validation}, and a sensitivity version of the cleaned dataset is exported with those ambiguous cases removed.

The output of Step~1 is therefore not only a cleaned occurrence table, but also an auditable record of what was standardized, what remained exact, what required conservative repair, and which cases still need manual confirmation before they are treated as fully trusted provenance.

\subsubsection{Outputs of this Step}
\label{sec:step1_outputs}

The corrected Step~1 produces the following files:
\begin{itemize}
    \item \texttt{01\_skills\_occurrence\_clean.json}
    \item \texttt{01\_skills\_occurrence\_clean.csv}
    \item \texttt{01\_skills\_occurrence\_clean\_sensitivity\_excluding\_ambiguous.json}
    \item \texttt{01\_skills\_occurrence\_clean\_sensitivity\_excluding\_ambiguous.csv}
    \item \texttt{01\_unique\_textual\_instances\_vectors.csv}
    \item \texttt{01\_data\_audit\_summary.csv}
    \item \texttt{01\_schema\_validation\_summary.csv}
    \item \texttt{01\_vector\_integrity\_report.csv}
    \item \texttt{01\_vector\_integrity\_summary.csv}
    \item \texttt{01\_duplicate\_report.csv}
    \item \texttt{01\_coverage\_by\_curricular\_origin.csv}
    \item \texttt{01\_metadata\_standardization\_log.csv}
    \item \texttt{01\_preprocessing\_decision\_log.csv}
    \item \texttt{01\_provenance\_resolution\_audit.csv}
    \item \texttt{01\_step1\_deviations\_notes.txt}
\end{itemize}

The cleaned occurrence dataset now includes the fields \texttt{text\_instance\_id}, \texttt{provenance\_resolution\_confidence}, and \texttt{requires\_manual\_validation}. The file \texttt{01\_unique\_textual\_instances\_vectors.csv} is the corrected screening base used in Step~2.

\subsubsection{Fields to Fill After Concluding This Step}
\label{sec:step1_fields_to_fill}

No extra analyst input is required to complete the computational part of Step~1. The only remaining manual task is to inspect the flagged provenance repairs recorded in:
\begin{itemize}
    \item \texttt{01\_preprocessing\_decision\_log.csv}
    \item \texttt{01\_provenance\_resolution\_audit.csv}
\end{itemize}

These files identify the raw records whose occurrence expansion depends on fallback provenance reconstruction and should therefore be checked before the fully trusted node universe is frozen.

\subsubsection{Implementation Note and Current Run}
\label{sec:step1_implementation_note}

The current corrected run produced:
\begin{itemize}
    \item 895 raw records;
    \item 1,038 occurrence rows after preprocessing;
    \item 26 raw records with list-length mismatch that required chunk-based provenance repair;
    \item 88 occurrence rows flagged with \texttt{requires\_manual\_validation = true};
    \item 950 rows in the sensitivity dataset that excludes those ambiguous cases;
    \item 895 unique textual instances in \texttt{01\_unique\_textual\_instances\_vectors.csv};
    \item 0 invalid \texttt{name\_vector} rows and 0 invalid \texttt{description\_vector} rows;
    \item 0 exact duplicate occurrence groups.
\end{itemize}

The main documented implementation choices are:
\begin{itemize}
    \item direct metadata broadcasting is preferred, with chunk-based reconstruction used only as a fallback;
    \item duplicate detection is exact-only at the occurrence level;
    \item normalization is conservative and does not rewrite the semantic content of the extracted skills;
    \item array-valued fields are serialized as JSON strings in CSV exports for traceability.
\end{itemize}
